\chapter{Cleft lip}
\index{Cleft lip}

\section*{Definition and Overview}
A cleft lip is a common congenital craniofacial anomaly characterised by an opening or split in the upper lip. This condition occurs during early foetal development when the tissues that form the lip fail to fuse completely. A cleft lip can range in severity from a minor notch in the vermilion border to a complete separation extending into the nostril. It can present unilaterally (on one side, more commonly the left) or bilaterally (on both sides), and may occur as an isolated defect or in combination with a cleft palate (cleft lip and palate). The condition has significant clinical implications, affecting feeding, speech development, dentition, and psychosocial well-being. It requires comprehensive, coordinated multidisciplinary care from infancy through adolescence.

\section*{Epidemiology}
Cleft lip, with or without cleft palate, is one of the most common congenital birth defects globally, occurring in approximately 1 in 700 to 1 in 1,000 live births. The prevalence varies significantly by geographic region and ethnic background, with the highest rates observed in Asian and Native American populations (up to 1 in 500 births) and the lowest in African populations (around 1 in 2,500 births). In many countries, the incidence is approximately 1 in 800 live births. Cleft lip is twice as common in males as in females. Approximately 70\% of cleft lip cases are non-syndromic, occurring in isolation, while the remaining 30\% are syndromic, associated with over 300 genetic syndromes, such as Van der Woude syndrome, Trisomy 13 (Patau syndrome), and Pierre Robin sequence.

\section*{Aetiology and Pathophysiology}
The development of a cleft lip involves a complex interplay of genetic susceptibility and environmental exposures during early embryogenesis.
\begin{itemize}
    \item \textbf{Embryological fusion failure:} Pathophysiologically, a cleft lip results from the failure of fusion between the medial nasal process and the maxillary processes during the fourth to seventh weeks of gestation.
    \item \textbf{Genetic mutations:} Multiple susceptibility genes have been identified, including interferon regulatory factor 6 ($IRF6$), ventral anterior homeobox 1 ($VAX1$), and ventral homeobox genes. A family history of clefting significantly increases recurrence risk.
    \item \textbf{Maternal smoking and alcohol exposure:} Prenatal exposure to tobacco smoke and alcohol significantly increases the risk of cleft development by inducing hypoxia and impairing cell migration in the developing face.
    \item \textbf{Nutritional deficiencies:} Lack of maternal folate (vitamin B9) intake before and during early pregnancy is a well-established risk factor, as folate is critical for DNA synthesis and tissue fusion.
    \item \textbf{Teratogenic medications:} Maternal use of certain anticonvulsants (e.g., phenytoin, valproic acid), retinoids, or corticosteroids during the first trimester is associated with a higher incidence of craniofacial clefts.
    \item \textbf{Maternal obesity and diabetes:} Pre-existing maternal diabetes or gestational metabolic disturbances can disrupt normal embryogenesis and increase malformation rates.
\end{itemize}

\section*{Clinical Features}
A cleft lip is visually apparent at birth and may also be detected prenatally via ultrasound.
\begin{itemize}
    \item \textbf{Visible split in the upper lip:} A gap that may be unilateral or bilateral, extending from the red border of the lip (vermilion) up towards or into the nostril.
    \item \textbf{Nasal deformity:} Flattish, widened nostril on the affected side due to the disruption of the orbicularis oris muscle insertion and nasal cartilage displacement.
    \item \textbf{Dental anomalies:} Misaligned, missing, or malformed teeth in the line of the cleft, particularly affecting the lateral incisors.
    \item \textbf{Feeding difficulties:} Difficulty establishing a seal around the nipple during breastfeeding or bottle-feeding, leading to air ingestion and poor weight gain, although this is less severe than in cleft palate.
    \item \textbf{Psychosocial challenges:} Parental distress at birth, and potential social, behavioural, or self-esteem challenges for the child as they grow.
\end{itemize}

\section*{Differential Diagnosis}
When a cleft lip is identified, clinicians must differentiate between non-syndromic (isolated) clefting and syndromic clefting associated with broader systemic anomalies.
\begin{itemize}
    \item \textbf{Non-syndromic cleft lip:} An isolated craniofacial defect with no associated developmental, intellectual, or systemic abnormalities.
    \item \textbf{Van der Woude syndrome:} An autosomal dominant disorder characterised by cleft lip and/or cleft palate accompanied by distinctive congenital pits or sinuses in the lower lip.
    \item \textbf{Trisomy 13 (Patau syndrome):} A severe chromosomal disorder characterised by midline cleft lip/palate, microphthalmia, polydactyly, and severe intellectual disability.
    \item \textbf{Amniotic band syndrome:} A non-genetic condition where fibrous bands of the amniotic sac constrict or shear embryonic tissue, causing atypical, asymmetric facial clefts that do not follow standard anatomical fusion lines.
    \item \textbf{EEC syndrome:} A genetic syndrome presenting with cleft lip/palate, ectodermal dysplasia (abnormal hair, nails, sweat glands), and split hand/foot malformations (ectrodactyly).
\end{itemize}

\section*{When to Seek Medical Care}
Parents of an infant born with a cleft lip should receive immediate support and referral to a cleft coordinator or specialised cleft palate team.

\textbf{Contact your local emergency services for:}
\begin{itemize}
    \item Signs of severe respiratory distress in the newborn, such as gasping, grunting, chest retractions, or cyanosis (blue-coloured skin/lips).
    \item Choking or severe aspiration during feeding, leading to persistent coughing or inability to breathe.
    \item Unresponsiveness, lethargy, or loss of consciousness in the infant.
\end{itemize}

\section*{Diagnosis and Investigations}
Diagnosis is made prenatally or immediately at birth, followed by comprehensive developmental evaluation.
\begin{itemize}
    \item \textbf{Prenatal ultrasound:} Often detected during the routine mid-trimester morphology scan (around 18--20 weeks gestation) using 2D, 3D, or 4D ultrasound imaging.
    \item \textbf{Clinical examination at birth:} Visual and physical inspection of the lip, palate, and nose, combined with a comprehensive neonatal examination to check for other congenital anomalies.
    \item \textbf{Genetic assessment and karyotyping:} Recommended if syndromic clefting is suspected due to the presence of other structural abnormalities or developmental delays.
    \item \textbf{Feeding assessment:} Conducted by a speech pathologist or lactation consultant to evaluate sucking efficiency, coordinate swallowing, and recommend specialized bottles.
    \item \textbf{Audiological screening:} Hearing tests (such as automated auditory brainstem response) to screen for middle ear effusion, which is common in children with clefts.
\end{itemize}

\section*{Management}
Management is multidisciplinary and extends from birth through to early adulthood, with surgical repair as the cornerstone of treatment.
\begin{itemize}
    \item \textbf{Multidisciplinary cleft team:} Collaboration between plastic surgeons, paediatricians, speech pathologists, orthodontists, paediatric dentists, ENT specialists, audiologists, geneticists, and psychologists.
    \item \textbf{Specialist feeding support:} Using specialized soft bottles and teats (e.g., Haberman feeder) to assist infants with feeding difficulties.
    \item \textbf{Primary surgical repair:} Surgical reconstruction of the lip, typically performed between 3 and 6 months of age, following the "Rule of 10s" (10 weeks old, 10 pounds in weight, and haemoglobin of 10\,g/dL).
    \item \textbf{Presurgical infant orthopaedics:} A non-surgical method (such as nasoalveolar moulding) using a custom moulding plate to approximate the lip segments and reshape the nose prior to surgery.
    \item \textbf{Speech therapy and dental monitoring:} Continuous support to monitor speech development and manage dental arch alignment as primary and secondary teeth erupt.
    \item \textbf{Secondary surgical revisions:} Nose and lip revisions (rhinoplasty, scar revision) performed during childhood or adolescence to improve symmetry and function.
\end{itemize}

\section*{Complications}
\begin{itemize}
    \item \textbf{Feeding insufficiency and poor weight gain:} Inability to feed effectively in the neonatal period, potentially leading to failure to thrive.
    \item \textbf{Dental malocclusion and crowded teeth:} Misalignment of the jaw and teeth in the cleft region, requiring extensive orthodontic treatment.
    \item \textbf{Speech articulation difficulties:} Difficulty producing certain speech sounds, particularly if there is co-existing cleft palate or velopharyngeal insufficiency.
    \item \textbf{Recurrent middle ear infections:} Increased risk of otitis media with effusion (glue ear) due to Eustachian tube dysfunction, which can cause hearing loss.
    \item \textbf{Psychosocial distress:} Self-esteem issues, bullying, or social anxiety during childhood and adolescence.
\end{itemize}

\section*{Prognosis and Outcomes}
The long-term prognosis for children with a cleft lip is excellent when managed by a dedicated multidisciplinary cleft team. Modern surgical techniques achieve outstanding aesthetic and functional outcomes, with minimal scarring and normal lip mobility. Most children grow up to lead healthy, active lives, with normal speech, hearing, and cognitive development. Ongoing orthodontic care and minor surgical adjustments may be required during childhood and adolescence, but final outcomes are highly successful.

\section*{Prevention and Self-Care}
\begin{itemize}
    \item \textbf{Folic acid supplementation:} Women planning a pregnancy should take a daily folic acid supplement (at least 400\,microg, or higher if at high risk) starting at least one month before conception and throughout the first trimester.
    \item \textbf{Avoid tobacco and alcohol:} Complete abstinence from smoking, vaping, and alcohol consumption during pregnancy to reduce the risk of clefting.
    \item \textbf{Review medications pre-conception:} Consult a GP or obstetrician before conceiving to ensure any chronic medications (e.g., anticonvulsants) are managed safely.
    \item \textbf{Manage pre-existing health conditions:} Optimise control of maternal diabetes and obesity before and during pregnancy.
    \item \textbf{Engage with support networks:} Reach out to organisations like CleftPALS (Cleft Palate and Lip Society) for peer support and educational resources.
\end{itemize}

\section*{Sources and Further Reading}
The following reputable organisations provide further information on cleft lip:
\begin{itemize}
    \item CleftPALS (Cleft Palate and Lip Society Australia). \textit{Support network and resources for families affected by cleft lip and palate.} https://www.cleftpals.org.au/
    \item Royal Children's Hospital Melbourne. \textit{Clinical guidelines and parent information on cleft lip repair.} https://www.rch.org.au/
    \item Sydney Children's Hospitals Network. \textit{Cleft lip and palate service information and treatment pathways.} https://www.schn.health.nsw.gov.au/
    \item Mayo Clinic. \textit{Overview of causes, diagnosis, and surgical management of cleft lip.} https://www.mayoclinic.org/
    \item National Health Service (NHS). \textit{Patient guide to cleft lip and palate care and support.} https://www.nhs.uk/
\end{itemize}
